\tocsection{\textbf{9/28/26 Lecture 14:} More on Topological Groups}

\begin{definition}[\textbf{Definition 14.1:} Left and Right Multiplication Maps]
    For a topological group $G$ and $x \in G$, define $L_{x} : G \to G$ by
    \[ L_{x}(g) = m(x, g) = x \cdot g \]
    and similarly, define $R_{x} : G \to G$ by
    \[ R_{x}(g) = m(g, x) = g \cdot x \]
\end{definition}

\begin{lemma}[\textbf{Lemma 14.2}]
    The maps $L_{x}$ and $R_{x}$ defined above are homeomorphisms.
\end{lemma}
\begin{proof}{Proof:}
    First, we will show that $L_{x}$ is continuous.\\
    
    Define $\phi : G \to G \times G$ by $\phi(g) = (x, g)$. If $U$ and $V$ are open in $G$, then $U \times V$ is a basis element for $G \times G$. Then, we see that
    \[ \phi^{-1}(U \times V) = \left\{\begin{matrix} \emptyset & \text{ if } x \notin U \\ V & \tt{ if } x \in U\end{matrix}\right. \]
    which is open. Thus, $\phi$ is continuous. Now, notice that $L_{x} = m \circ \phi$, so thus, $L_{x}$ is continuous as $m$ is.\\

    Then, as $(L_{x})^{-1} = L_{x^{-1}}$, we see that $L_{x}$ has a continuous inverse, so it is a homeomorphism. The same proof works for $R_{x}$. 
\end{proof}

\begin{lemma}[\textbf{Corollary 14.3}]
    For a topological group $G$ and any $g, g' \in G$, there exists a homeomorphism $f : G \to G$ for which $f(g) = g'$. 
\end{lemma}
\begin{proof}{Proof:}
    Let $f = L_{g'g^{-1}}$. This is a homeomorphism satisfying the above property. 
\end{proof}

\begin{definition}[\textbf{Definition 14.4:} Homogeneous Space]
    A topological space $X$ is \emph{homogeneous} if, for any pair of points $x, y \in X$, there is a homeomorphism $f : X \to X$ such that $f(x) = y$. 
\end{definition}

\begin{lemma}[\textbf{Proposition 14.5}]
    If $G$ is a topological group, and $K$ is the connected component containing $e$, then $K$ is a closed normal subgroup of $G$. 
\end{lemma}
\begin{proof}{Proof:}
    First, we see that $K$ is closed as all components are closed. Now, we need to show that $K$ is a normal subgroup of $G$.\\

    Notice that for any $x \in K$, $Kx^{-1} = R_{x^{-1}}(K)$, so $Kx^{-1}$ is connected. However, since $x \in K$, $e = xx^{-1} \in R_{x^{-1}}(K)$, so $R_{x^{-1}}(K)$ is the connected component containing $e$, which is contained in $K$. Then, as $e \in K$, we see that $x^{-1} \in R_{x^{-1}}(K)$, so $x^{-1} \in K$. Thus, $K$ is closed under inverses.\\

    Now, if $y \in K$, then $yx^{-1} \in K$, as $yx^{-1} \in R_{x^{-1}}(K)$, which implies that $K$ is closed under products.\\

    Finally, we need to show that $K$ is a normal subgroup. Let $g \in G$. Then, we see that
    \[ gKg^{-1} = L_{g}(R_{g^{-1}}(K)) \]
    This is connected as $L_{g}$ and $R_{g^{-1}}$ are continuous, but also $e \in gKg^{-1}$ as $e \in K$. Thus, as $gKg^{-1}$ is a connected component containing $e$, it must be in $K$. Thus, $gKg^{-1} \subseteq K$, so $K$ is normal.

    $\qed$
\end{proof}

\begin{lemma}[\textbf{Proposition 14.6}]
    Suppose that $G$ is a connected topological group, and $V$ is an open neighborhood of $e$. Then, $G$ is generated by $V$. 
\end{lemma}
\begin{proof}{Proof:}
    Let $H = \brac{V}$ be the subgroup generated by $V$. Then, for all $h \in H$, we can write
    \[ h = v_1^{n_1}v_2^{n_2} \cdots v_{k}^{n_k} \]
    with $v_i \in V$ and $n_i \in \Z$. We want to show that $H$ is both open and closed. To see that $H$ is open, let $h \in H$ and consider $hV$. We have that $hV \subseteq H$ by definition of $H$. Well, $hV = L_{h}(V)$, so as $L_{h}$ is a homeomorphism, $L_{h}(V)$ is an open neighborhood of $h$ ($h \in hV$ as $e \in V$).\\

    Now, to show that $H$ is closed. We want to show that $G \setminus H$ is open. Let $g \in G \setminus H$. Now, $gV = L_{g}(V)$ is an open neighborhood of $g$ by above, so we want to show that $gV \subseteq G \setminus H$. To do this, notice that if $x \in H$ and $x \in gV$, then we can write $x = gv$ for some $v \in V$, but when we see that $g = xv^{-1}$, and as $xv^{-1} \in H$, this implies that $g \in H$, which is a contradiction. Thus, we must have that $gV \subseteq G \setminus H$, so $G \setminus H$ is open.\\

    Thus, as $G$ is connected, either $H = G$ or $H = \emptyset$, but $H$ contains $e$, so it must be the former. 

    $\qed$
\end{proof}

\begin{lemma}[\textbf{Proposition 14.7}]
    Consider the group $\tt{GL}(n) = \{\tt{invertible } n \times n \tt{ real matrices}\}$ with the topology coming from the homeomorphism with $\R^{n^2}$. Then, $\tt{GL}(n)$ is a topological group.
\end{lemma}
\begin{proof}{Proof:}
    First, we want to show that $m : M_{n} \times M_{n} \to M_{n}$ is continuous. Well, we have that
    \[ (AB)_{ij} = \sum_{k = 1}^{n} a_{ik}b_{kj} \]
    is a polynomial in $a_{ij}$ and $b_{ij}$, which are continuous on $\R^2$. We also have that
    \[ (A^{-1})_{ij} = \frac{1}{\det(A)} (ji^{th} \tt{ cofactor of A}) \]
    Then, as this is a rational function in the entires of $(a_{ij})$ with a determinant that is never $0$, we see that inversion is also continuous.

    $\qed$
\end{proof}

\begin{lemma}[\textbf{Proposition 14.8}]
    Consider the following subgroups of $\tt{GL}(n)$:
    \[ O(n) = \{A \in \tt{GL}(n) \col A^{-1} = A^{T}\} \]
    and
    \[ SO(n) = \{A \in O(n) \col \det(A) = 1\} \]
    Then, $O(n)$ and $SO(n)$ are compact. 
\end{lemma}
\begin{proof}{Proof:}
    Start with $O(n)$. To show that it is closed and bounded in $M_n$, as this is isomorphic to $\R^{n^2}$. Define $\phi : M_{n} \to M_{n}$ by $\phi(M) = MM^{T}$. Then, we see that
    \[ O(n) = \phi^{-1}(I_{n}) \]
    so as $\phi$ is continuous, and as $\{I_{n}\}$ is closed, we see that $O(n)$ is closed. Now, to show that it is bounded, notice that
    \[ A \in O(n) \implies \tt{tr}(AA^{T}) = \tt{tr}(I) \implies \sum_{i, j = 1}^{n} a_{ij}^2 = n \]
    thus, we see that $O(n)$ is bounded in $M_{n}$.\\

    Now, we will do $SO(n)$. First, notice that $\det : O(n) \to \R$ is continuous, so as $SO(n) = \det^{-1}(1)$, and $\{1\}$ is closed, we see that $SO(n)$ is closed in $O(n)$, so it is compact. 

    $\qed$
\end{proof}

Here are some examples of these groups and what spaces they are similar to.
\begin{itemize}
    \item $n = 1$: We see that $SO(n) = \{1\}$ and $O(1) = \{1, -1\}$\\
    \item $n = 2$: We have
    \[ SO(2) = \left\{\begin{pmatrix} \cos \theta & -\sin \theta \\ \sin \theta & \cos \theta\end{pmatrix}\right\} = S^1 \]
    and $O(2)$ is homeomorphic to a disjoint union of two circles; $SO(2)$ and $SO(2)$ times the matrix
    \[ \begin{pmatrix} 1 & 0 \\ 0 & -1 \end{pmatrix} \]
    \item $n = 3$: We have that $SO(3)$ is homeomorphic to $P^3$, three dimensional projective space. 
\end{itemize}